\chapter{Executive Summary}

This project designs and validates a collaborative document editor with an integrated AI writing
assistant. The target product is intentionally framed as a simplified Google Docs competitor: multiple
users can open the same document, see each other's edits in real time, manage sharing and access,
invoke AI assistance on selected text, inspect document history, and export the result. The assignment
brief asks not only for a prototype, but for evidence of disciplined requirements engineering,
architecture, repository organization, and team coordination.

The final system follows a modular monorepo structure with distinct frontend, backend, realtime,
AI-service, shared, and documentation layers. The frontend is a React application focused on document
editing, collaboration, permissions, AI suggestions, version history, and export. The backend is an
Express API responsible for authentication, document lifecycle management, role-based access control,
AI orchestration, export jobs, versioning, and realtime session issuance. The realtime service is a
dedicated WebSocket/Yjs relay that enforces collaboration permissions and propagates awareness and
content updates. Persistence is handled through SQLite for portability and deterministic setup, while
AI generation is abstracted behind a provider interface supporting both a deterministic stub provider
and a local LM Studio/OpenAI-compatible endpoint.

Several architectural drivers had the strongest influence on the design. First, collaborative editing
had to feel responsive and correct; this drove the decision to use Yjs-based push synchronization for
document state instead of periodic polling. Second, AI needed to be useful but failure-tolerant; this
led to an asynchronous job model with explicit \texttt{PENDING}/\texttt{RUNNING}/\texttt{SUCCEEDED}/\texttt{FAILED}
states rather than blocking inline model calls. Third, the codebase had to be maintainable by a small
student team; this motivated a monorepo, shared contract notes, explicit service boundaries, and
root-level orchestration scripts for installation, running, testing, and demo reset.

The delivered MVP intentionally goes beyond the assignment's minimum proof-of-concept threshold.
The original PoC requirement was a working frontend communicating with a backend over at least one
meaningful API, with data contracts matching the design. The final implementation still satisfies that
minimum, but also demonstrates live collaboration, autosave, permission management, version history,
revert, AI rewrite/summarize/translate jobs, partial AI acceptance, export jobs, Swagger/OpenAPI
docs, CI automation, and Docker Compose configuration. At the same time, this report is explicit about
MVP simplifications: authentication is local bearer-token login rather than a full external identity
provider, the storage engine is SQLite rather than a production-scale RDBMS deployment, and the editor
focuses on collaborative text rather than a full structured rich-text schema.

The report is therefore structured to do two things at once: satisfy the assignment rubric as a
requirements-and-architecture deliverable, and explain how the implemented MVP maps back to those
design decisions. Each major architecture section includes a delivered-MVP note where the practical
implementation intentionally differs from the longer-term target architecture. This keeps the report
truthful, versionable, and aligned with the actual repository and demo.
